% SPDX-FileCopyrightText: 2026 Will Cook
% SPDX-License-Identifier: CC-BY-4.0
%
% PROVENANCE: superseded historical draft.  The published successor is
% paper/erdos-251-prime-gap-dyadic-series.tex in the public release repository
% wcook04/plectis-lean-erdos249-257 at 8df0dda280 (release/final-20260905).
% Every label in this file is also defined there; do not export or promote
% from this copy.
%
% One of the Erdős Problem Notes: a short standalone treatment of a single
% problem in the ErdosProblems expansion library.  Build with
% `tectonic erdos-251-prime-gap-dyadic-series.tex` or `make`.
\documentclass[11pt]{article}
\input{problem-note-preamble}
% Problem-specific source pin: #251 advances independently of the corpus
% default, while the #249/#257 notes retain their own reviewed snapshot.
\renewcommand{\commit}{6bad8cba6327a4ceb4a1179cf8176597fdf2e381}
\renewcommand{\commitshort}{6bad8cba6327}
\hypersetup{
  pdftitle={A countermodel for growth-and-parity arguments on the prime-gap dyadic series (Erdős Problem 251)},
  pdfsubject={Erdős Problem Notes: Problem 251},
  pdfkeywords={irrationality, prime gaps, dyadic series, summation by parts, Lean 4},
}

\runninghead{Erd\H{o}s \#251: a countermodel and a denominator floor}
\title{A Countermodel for Growth-and-Parity Arguments\\on the Prime-Gap Dyadic Series}
\subtitle{An explicit denominator floor, the exact shift and LCM-diagonal
coordinates, and the open boundary for Erd\H{o}s Problem~\#251}
\author{Will Cook}
\date{July 2026}

\begin{document}
\maketitle
\noteprovenance

\begin{abstract}
The positive even word $a_n=2(n^2+4n+2)$ has complete dyadic tails
$U_N=2(N+4)^2$ and rational sum $\sum_{j\ge1}a_j2^{-j}=32$.
Its unbounded, nonperiodic coefficients show why growth and parity do not
force the small-tail configurations used in a natural irrationality argument.
A complementary construction preserves any finite prime-gap prefix and
residues modulo one prescribed integer while making the dyadic sum rational
by a bounded upward perturbation.  We then identify the complete tails of
the actual prime-gap series and classify rationality by their integral
shifts.  Two adjacent small shifts with unequal associated gaps give a local
obstruction to integrality.  Supplying those configurations cofinally for
every positive shift would prove irrationality of the prime series; that
prime-specific assertion remains unproved.
\end{abstract}

\section{Introduction}\label{sec:problem}

\paragraph{The complete countermodel in one calculation.}
Set $a_n=2(n^2+4n+2)$ and $U_N=2(N+4)^2$.  Since
$2U_N-U_{N+1}=a_{N+1}$, finite telescoping gives
\[
 \sum_{j=1}^{m}\frac{a_{N+j}}{2^j}
   =U_N-\frac{U_{N+m}}{2^m}\longrightarrow U_N.
\]
Thus $U_N$ is the actual infinite tail, not merely a solution of its
recurrence, and the sum at $N=0$ is $32$.  Moreover
\[
 a_{n+1}-a_n=4n+10,\qquad
 U_{N+h}-U_N=2h(2N+h+8)\in\Z.
\]
These positive even coefficients grow polynomially and are unbounded and
nonperiodic, but every tail difference is integral.  The unshifted
normalisation $\sum_{n\ge0}a_n2^{-(n+1)}$ is $18$.
Proposition~\ref{res:polynomialcountermodel} states the resulting obstruction.

\paragraph{Preserving more of the actual gaps.}
The quadratic example has the wrong cumulative growth to imitate the primes.
A different construction addresses that limitation.  If
$A=\sum_{n\ge0}a_n2^{-(n+1)}$ converges, fix $M\ge1$ and a cutoff $K$, and
choose a rational $r\in(A,A+M2^{-K})$.  Write
\[
 \frac{2^K(r-A)}{M}=\sum_{j\ge0}\frac{\delta_j}{2^{j+1}},
 \qquad \delta_j\in\{0,1\}.
\]
Set $b_n=a_n$ for $n<K$ and $b_{K+j}=a_{K+j}+M\delta_j$.  Then
$\sum_{n\ge0}b_n2^{-(n+1)}=r$.  Every coefficient retains its residue
modulo $M$ and changes by at most $M$; each cumulative sum changes by at
most $Mn$.  Applied to prime gaps, this retains the cumulative
$n\log n$ scale.  Taking $M$ even retains parity as well.  Thus a bounded
residue-preserving perturbation can have a rational value without sharing
the polynomial example's growth.  Neither construction asserts that its
cumulative sums are prime.

\paragraph{A stronger obstruction.}
The bounded perturbation preserves one modulus.  Proposition~\ref{res:sparserationalisation}
uses variable pair capacities, a congruence buffer, and sites of upper Banach
density zero.  It realises an interval of dyadic values while preserving
every fixed modulus and the empirical laws of all unnormalised gap blocks
of length $o(\log\log X)$.  The pair/buffer algebra and variable-alphabet
filling are Lean-checked; the sparse location schedule is ordinary.
It does not prove a prime-gap distribution law, a circular-arc lower bound,
or irrationality of $\Pi$.  The first numbered theorem remains the
lcm-diagonal classifier of the actual tails.

\paragraph{The prime series and its tails.}
Index the primes from zero: $p_0=2$, $p_1=3$, and
$g_n=p_{n+1}-p_n$.  In this notation Erd\H{o}s Problem~\#251 asks whether
\[
 \Pi=\sum_{n\ge0}\frac{p_n}{2^{n+1}}
 \quad\text{is irrational}.
\]
This is the usual one-based series $2/2+3/4+5/8+\cdots$.
The first gaps are $1,2,2,4,2,4,2,4,6,2$; only $g_0$ is odd.
Let $G=\sum_{n\ge0}g_n2^{-(n+1)}$.  The exact identities are
\[
 \Pi=2+G,\qquad
 T_N=\sum_{j\ge1}\frac{g_{N+j}}{2^j},\qquad
 T_0=2G-1,\qquad T_{N+1}=2T_N-g_{N+1}.
\]
Convergence and the identification of these actual tails are established
below.  They are not additional assumptions under a hypothetical rational
value of $\Pi$.

\paragraph{What remains to be distinguished.}
For any integer-digit recurrence, iteration gives
\[
 T_N=2^NT_0-B_N,\qquad
 T_{N+h}-T_N=2^N(2^h-1)T_0-C_{N,h},
 \qquad B_N,C_{N,h}\in\Z.
\]
A rational initial value forces an eventual integral shift; an irrational
initial value permits no integral positive shift.  This is the arithmetic
reason for the lcm-diagonal criterion in Theorem~\ref{res:lcmdiagonal}.
A more local sufficient condition uses two adjacent shift values in
$(-1,1)$ and a nonzero associated gap difference.  If both values were
integral they would be zero, contradicting their recurrence.  The missing
input is their joint occurrence at arbitrarily late indices for every
positive shift.  The countermodels show why several coarse hypotheses do
not supply it.

\paragraph{Relation to prior work.}
Erd\H{o}s stated the fixed-denominator question in his work on irrational
series~\cite[p.~94]{erdos1958}\cite[p.~62]{erdosgraham1980}%
\cite[p.~103]{erdos1988}.  His irrationality result for prime numerators
with factorial denominators does not apply to this dyadic denominator
sequence.  Kova\v{c}'s counterexample to a related variable-denominator
expectation likewise does not decide the fixed-denominator problem
\cite{kovac2026}.  An adjacent dyadic theorem of Erd\H{o}s and Pomerance
concerns bounded digits recording comparisons of the largest prime factors
of consecutive integers~\cite[\S7]{erdospomerance1978}; our prime numerators
and gaps are unbounded.  The detailed historical comparisons and their
hypotheses belong to the companion reasoning record.  No priority claim
is made for the elementary identities or constructions in this note.

\paragraph{Notation and organisation.}
Write $\N=\{0,1,\ldots\}$, $\Z$ for the integers,
$\operatorname{Irr}(x)$ for irrationality, $\operatorname{den}(q)$ for the
reduced denominator of $q\in\Q$, and $\ph$ for Euler's totient.
Section~\ref{sec:diagonal-collapse} isolates the recurrence-level
classifiers and the polynomial obstruction.
Section~\ref{sec:parts} proves the actual prime-to-gap identity;
Section~\ref{sec:tail} proves the shift classification and local criterion.
Section~\ref{sec:carry} explains the remaining carry freedom, and
Section~\ref{sec:open} states the exact prime-specific obligation.
The sources and verification boundary are recorded after the mathematics.

\medskip
\noindent\textbf{Keywords.} irrationality; prime gaps; dyadic series;
summation by parts; Lean~4.
\qquad\textbf{MSC 2020.} 11J72 (primary); 11N05, 68V20 (secondary).

\section{One diagonal, and two circular escape tests}\label{sec:diagonal-collapse}

We first record the strongest recurrence-level conclusions.  They apply to
arbitrary integer digits and therefore separate the finite algebra from the
prime-specific input still missing in Problem~\#251.

For a real recurrence put
\[
 \Delta_hT(N)=T_{N+h}-T_N,
\]
and let $L_0=1$, $L_j=\operatorname{lcm}(1,\ldots,j)$ for $j\ge1$.

\begin{theorem}[lcm-diagonal criterion]\label{res:lcmdiagonal}
Let $g:\N\to\Z$ and $T:\N\to\R$ satisfy
$T_{N+1}=2T_N-g_{N+1}$ for every $N$.  Then
\[
 \operatorname{Irr}(T_0)
 \quad\Longleftrightarrow\quad
 \Delta_{L_j}T(L_j)\notin\Z
 \quad\hbox{for every }j\ge0.
\]
\end{theorem}

\begin{proof}
The exact rationality classification proved in
Theorem~\ref{res:escape-irrational} says that an irrational initial state has
no integral positive shift at any basepoint.  This gives the forward
implication.  Conversely, if $T_0$ is rational, some positive shift length
$h$ is integral at every basepoint beyond an index $N_0$.  Choose $j$ so large
that $N_0\le L_j$ and $h\mid L_j$.  The cocycle identity
\[
 \Delta_{a+b}T(N)=\Delta_aT(N)+\Delta_bT(N+a)
\]
shows by induction that every positive multiple of $h$ is integral at every
such basepoint.  Hence $\Delta_{L_j}T(L_j)$ is integral, contradicting the
right-hand side.
\end{proof}

The factorial schedule $j!$ has the same divisibility property, but the lcm
schedule is pointwise no larger and is the canonical endpoint used here.  The
theorem is an exact one-sequence reformulation; it gives no non-integrality
statement for the actual prime-gap recurrence.

The adjacent-small-mismatch condition also admits an exact normal form.  For a
rational recurrence and a fixed $h$, write
\[
 D_N=\Delta_hT(N),\qquad
 \delta_N=g_{N+h+1}-g_{N+1}.
\]
Then $D_{N+1}=2D_N-\delta_N$.

\begin{theorem}[signed two-window normal form]\label{res:signedwindow}
Assume that $\delta_N$ is even.  The conjunction
\[
 -1<D_N<1,\qquad -1<D_{N+1}<1,\qquad \delta_N\ne0
\]
is equivalent to
\[
 \bigl(\delta_N=2\ \hbox{ and }\tfrac12<D_N<1\bigr)
 \quad\hbox{or}\quad
 \bigl(\delta_N=-2\ \hbox{ and }-1<D_N<-\tfrac12\bigr).
\]
\end{theorem}

\begin{proof}
The recurrence and the two unit windows give $-3<\delta_N<3$.  A nonzero even
integer in that interval is $2$ or $-2$.  Substitution into
$D_{N+1}=2D_N-\delta_N$ gives the stated half-window and, in the reverse
direction, recovers the second unit window.
\end{proof}

For the longer block, set
\[
 B_{h,N,r}=\sum_{i=0}^{r-1}2^{r-1-i}\delta_{N+i};
 \qquad D_{N+r}=2^rD_N-B_{h,N,r}.
\]
Call the data-dependent affine condition
\[
 \mathcal A_{N,r}:\qquad
 D_{N+r}\in -B_{h,N,r}+2^{r+1}\Z.
\]

\begin{theorem}[affine and fixed-lattice circularity]\label{res:affinecollapse}
For every rational dyadic tail recurrence and all $h,N,r\ge0$,
\[
 \mathcal A_{N,r}\quad\Longleftrightarrow\quad D_N\in2\Z.
\tag{2.1}\label{eq:affinecollapse}
\]
Consequently, if every $\delta_N$ is even, then
\[
 \bigl(\forall N_0\ \exists N,r:\ N_0<N\text{ and }\neg\mathcal A_{N,r}\bigr)
 \quad\Longleftrightarrow\quad
 D_N\notin\Z\text{ for arbitrarily large }N.
\tag{2.2}\label{eq:affinecofinal}
\]

There is a second equivalence.  Let $b:\N\to\Q$ satisfy
$|D_N|\le b(N)$ for every $N$, and suppose that for every $N$ and every
positive integer $q$ there is an $r$ with
\[
 2b(N+r)q<2^r.
\tag{2.3}\label{eq:dyadicscale}
\]
Then
\[
 \begin{split}
 &\forall N_0\ \exists N,r:\ N_0<N\text{ and }
   \forall z\in\Z,\quad
   b(N+r)<|B_{h,N,r}-2^rz|\\
 &\hspace{35mm}\Longleftrightarrow\quad
 D_N\notin\Z\text{ for arbitrarily large }N.
 \end{split}
\tag{2.4}\label{eq:fixedcollapse}
\]
\end{theorem}

\begin{proof}
Substituting $D_{N+r}=2^rD_N-B_{h,N,r}$ into $\mathcal A_{N,r}$ cancels the
observed block from both sides and leaves $D_N=2z$.  This proves
\eqref{eq:affinecollapse}.  After one recurrence step, evenness of
$\delta_N$ identifies even integrality at $N+1$ with ordinary integrality at
$N$.  The forward implication in~\eqref{eq:affinecofinal} follows, while the
reverse implication chooses a nonintegral $D_N$ and the legal depth $r=0$.

For~\eqref{eq:fixedcollapse}, eventual integrality and the block identity give
some $z\in\Z$ with
\[
 B_{h,N,r}-2^rz=-D_{N+r},
\]
so the displayed strict separation contradicts $|D_{N+r}|\le b(N+r)$.
Conversely, if $D_N$ is nonintegral and has reduced denominator $q$, then its
distance from every integer is at least $1/q$.  Choose $r$ from
\eqref{eq:dyadicscale}, scale this separation by $2^r$, and use the block
identity together with $|D_{N+r}|\le b(N+r)$.  The triangle inequality gives
$|B_{h,N,r}-2^rz|>b(N+r)$ for every integer $z$.
\end{proof}

The first equivalence shows that the apparent affine hierarchy contains no
depth-dependent information: it is the pullback of even integrality through
the recurrence identity.  The fixed lattice removes that data dependence, but
under the stated growth condition rational denominator separation already
forces every required escape.  Hence neither cofinal condition is an
independent source of information about consecutive primes.

Finally, the coarse coefficient profile itself has an exact infinite
counterexample.

\begin{proposition}[quadratic polynomial-shift countermodel]\label{res:polynomialcountermodel}
Put
\[
 a_n=2(n^2+4n+2),\qquad U_n=2(n+4)^2.
\]
Then $a_n$ is positive, even and strictly increasing,
$U_{n+1}=2U_n-a_{n+1}$, every shift $U_{N+h}-U_N$ is integral,
\[
 a_{N+2}-a_{N+1}\ne2,-2
 \qquad\hbox{for every }N,
\]
and
\[
 \sum_{j\ge1}\frac{a_j}{2^j}=32
\]
as the kernel-checked identity
\lword{Erdos251/PolynomialGapSeriesValue.lean}{92}
      {tsum_polynomialGapDyadicTerm_eq}{the countermodel series value}.
\end{proposition}

\begin{proof}
Expand $a_{n+1}-a_n=4n+10$ and $2U_n-U_{n+1}=a_{n+1}$.  The $h$-shift is
$U_{N+h}-U_N=2h(2N+h+8)$, an even integer.  The finite telescope
$\sum_{j=1}^{N}a_j/2^j=32-U_N/2^N$ is exact rational algebra from the
recurrence; $U_N/2^N\to0$ because a quadratic is dominated by $2^N$.  The
nonnegative terms therefore have infinite sum $32$.  Strict growth makes the
word unbounded and nonperiodic.
\end{proof}

\begin{proposition}[sparse congruence-preserving rationalisation]\label{res:sparserationalisation}
Let $a_n$ be nonnegative integers with $\sum_{n\ge0}a_n2^{-(n+1)}=A<\infty$,
let $K\in\N$, and let $f:\N\to\R$ tend to $+\infty$.  There exist a set
$S\subseteq[K,\infty)$ of upper Banach density zero and a nondegenerate
interval $I\subset(A,\infty)$ such that every $r\in I$ is realised by a
nonnegative correction $e$ supported on $S$, with $0\le e_n\le f(n)$
eventually, $\sum(a_n+e_n)2^{-(n+1)}=r$, and every fixed modulus eventually
dividing both $e_n$ and $\sum_{i<n}e_i$.  For the actual prime gaps and
every $0<\varepsilon<1$ the same construction may be taken with
$e_n\le(\log(n+3))^\varepsilon$ eventually and
$|S\cap[X,2X)|=O_\varepsilon(X/\log\log X)$, so that all unnormalised blocks
of length $o(\log\log X)$ are preserved in total variation.
\end{proposition}

\begin{proof}[Proof sketch]
Neighbouring corrections $M d$ and $M(D-d)$ have digit-independent unweighted
total $MD$ and one free weighted digit.  A buffer $c=(-C)\bmod M$ repairs the
new cumulative modulus before the digits are chosen and preserves earlier
moduli because $M_{j-1}\mid M_j$.  Capacities $D_j=2^{\ell_{j-1}}-1$ use the
preceding spacing, so the tail overlaps the current weight.  Greedy selection
then fills a nondegenerate interval.  The pair/buffer identities, the
capacity telescope, and variable-alphabet filling under those hypotheses are
the Lean declarations
\texttt{pair\_total}, \texttt{pair\_weighted}, \texttt{repairBuffer\_repairs},
\texttt{capacity\_telescope}, and \texttt{exists\_digits\_hasSum\_of\_capacity}
in \texttt{SparseRationalisationCore.lean}.  The location schedule, Banach
density, and block-law transfer are ordinary
(companion note \texttt{SparseRationalisation.md}).
This does not prove a prime-gap distribution law, a circular-arc lower bound,
or irrationality of $\Pi$.
\end{proof}

\section{Summation by parts, with the endpoint retained}\label{sec:parts}

Summation by parts trades a sequence for its consecutive differences.  The form
recorded here is exact: it carries no error term, it assumes nothing about the
sequence, and it retains the endpoint term rather than absorbing it into an
estimate.  For a sequence $P$ of rational numbers and $n\ge0$ put
\[
 D(P,n)=\sum_{i=0}^{n-1}\frac{P(i)}{2^{\,i+1}},
 \qquad
 \Delta(P,n)=\sum_{i=0}^{n-1}\frac{P(i+1)-P(i)}{2^{\,i+1}} ,
\]
both empty, hence zero, at $n=0$.  We call these the
\lword{Erdos251/PrimeGapDyadicTail.lean}{101}{dyadicPartialSumQ}{dyadic partial
sum} and the
\lword{Erdos251/PrimeGapDyadicTail.lean}{121}{dyadicDifferencePartialSumQ}{dyadic
difference sum} of $P$.

\begin{proposition}[finite summation by parts]\label{res:abel}
For every rational sequence $P$ and every $n\ge0$,
\[
 D(P,n+1)=P(0)+\Delta(P,n)-\frac{P(n)}{2^{\,n+1}} .
\]
\end{proposition}

\begin{proof}
A routine induction on $n$.  At $n=0$ both sides equal $P(0)/2$.  For the step, adding
$P(n+1)/2^{\,n+2}$ to the left and
$\bigl(P(n+1)-P(n)\bigr)/2^{\,n+1}$ to the difference sum changes the endpoint
term from $P(n)/2^{\,n+1}$ to $P(n+1)/2^{\,n+2}$, and the two adjustments agree.
\end{proof}

\noindent Formalised as the
\lword{Erdos251/PrimeGapDyadicTail.lean}{138}{dyadicPartialSumQ_eq_start_add_differences}{summation-by-parts
identity}.  Nothing is assumed about $P$: no positivity, no monotonicity, and no
convergence.

Specialising to $P(i)=p_i$, whose first value is $p_0=2$, and writing
$g_i=p_{i+1}-p_i$ for the zero-based gaps, gives the reformulation.

\begin{theorem}[prime-gap reformulation]\label{res:parts}
Let $p_0=2,p_1=3,\ldots$ be the primes in increasing order and
$g_i=p_{i+1}-p_i$.  For every $n\ge0$,
\[
 \sum_{i=0}^{n}\frac{p_i}{2^{\,i+1}}
 =2+\sum_{i=0}^{n-1}\frac{g_i}{2^{\,i+1}}-\frac{p_n}{2^{\,n+1}} .
\]
\end{theorem}

\noindent Formalised as the
\lword{Erdos251/PrimeGapDyadicTail.lean}{172}{prime0_dyadic_summation_by_parts}{prime-gap
summation by parts}, using the
\lword{Erdos251/PrimeGapDyadicTail.lean}{161}{primeGapPartialSumQ}{gap partial
sum} and the
\lword{Erdos251/PrimeGapDyadicTail.lean}{156}{primeGap0_cast}{gap cast
identity}; the latter records that the natural-number difference
$p_{n+1}-p_n$ agrees with the difference taken in $\Q$, which needs
$p_n\le p_{n+1}$, the
\lword{Erdos251/PrimeGapDyadicTail.lean}{152}{prime0_mono_step}{monotonicity
step}.  The leading $2$ is the first prime, not a normalising constant.
At $n=2$, for instance, the left side is $2/2+3/4+5/8=19/8$ and the right side
is $2+(1/2+2/4)-5/8=19/8$.

\subsection{The infinite identity and the irrationality equivalence}\label{sec:infinite}

Write
\[
 u_n=\frac{p_n}{2^{\,n+1}},\qquad
 v_n=\frac{g_n}{2^{\,n+1}}
\]
for the terms of the prime series and of the gap series, so that
$\sum_{n\ge0}u_n=\Pi$.  The termwise identity $v_n=2u_{n+1}-u_n$ is the
\lword{Erdos251/PrimeGapDyadicTail.lean}{202}{primeGapDyadicTerm_eq}{dyadic
discrete derivative}.  It expresses each gap term as an integer combination of
two consecutive prime terms, so once $(u_n)$ is summable the gap series can be
summed by rearranging two copies of the prime series, which is what the
following proof does.

\begin{theorem}[infinite prime-gap identity]\label{res:infinite}
Let $u_n=p_n/2^{\,n+1}$ and $v_n=g_n/2^{\,n+1}$.  If $(u_n)$ is summable, then
$(v_n)$ is summable and
\[
 \sum_{n\ge0}u_n=2+\sum_{n\ge0}v_n .
\]
\end{theorem}

\begin{proof}
The shifted sequence $(u_{n+1})$ is summable.  Sum
$v_n=2u_{n+1}-u_n$ and use $u_0=1$:
\[
 \sum_{n\ge0}v_n
 =2\left(\sum_{n\ge0}u_n-u_0\right)-\sum_{n\ge0}u_n
 =\sum_{n\ge0}u_n-2 .
\]
\end{proof}

\noindent The summability transfer is the
\lword{Erdos251/PrimeGapDyadicTail.lean}{385}{summable_primeGapDyadicTerm_of_summable_primeDyadicTerm}{gap-series
summability theorem}, and the displayed identity is the
\lword{Erdos251/PrimeGapDyadicTail.lean}{404}{tsum_primeDyadicTerm_eq_two_add_primeGap}{infinite
prime-gap identity}.

\begin{corollary}[exact irrationality reformulation]\label{res:irr-equivalence}
With $u_n$ and $v_n$ as in Theorem~\ref{res:infinite}, and $(u_n)$ summable,
\[
 \operatorname{Irr}\!\left(\sum_{n\ge0}u_n\right)
 \quad\Longleftrightarrow\quad
 \operatorname{Irr}\!\left(\sum_{n\ge0}v_n\right).
\]
The corresponding zero-based series with denominator $2^n$ is
$4+2\sum_{n\ge0}v_n$ and has the same irrationality status.
\end{corollary}

\noindent These are the
\lword{Erdos251/PrimeGapDyadicTail.lean}{435}{irrational_tsum_primeDyadicTerm_iff_primeGap}{normalised
irrationality equivalence}, the
\lword{Erdos251/PrimeGapDyadicTail.lean}{444}{tsum_primeDisplayedDyadicTerm_eq_four_add_two_primeGap}{displayed-series
identity}, and the
\lword{Erdos251/PrimeGapDyadicTail.lean}{459}{irrational_tsum_primeDisplayedDyadicTerm_iff_primeGap}{displayed-series
irrationality equivalence}.

\paragraph{The summability hypothesis.}
The formal source now proves the elementary polynomial bound
$p_n\le1250(n+1)^4$, using prime counting and central-binomial growth, and
deduces summability directly
(\lword{Erdos251/PrimeGapDyadicTail.lean}{360}{prime0_le_polynomial}{polynomial bound},
\lword{Erdos251/PrimeGapDyadicTail.lean}{379}{summable_primeDyadicTerm}{summability}).
Thus Theorem~\ref{res:infinite} and Corollary~\ref{res:irr-equivalence} are now
unconditional Lean-checked statements; the prime number theorem remains useful
context but is no longer a proof dependency of this note.  The open content is
exactly irrationality of the prime-gap series.  Concretely,
$\Pi=3.674643966\ldots$ is irrational if and only if
$\sum_{n\ge0}g_n2^{-(n+1)}=1.674643966\ldots$ is, and neither is known.

\section{The tail recurrence and integral shifts}\label{sec:tail}

The series is not attacked directly.  Suppose $\sum_{i\ge0}a_i2^{-(i+1)}$
converges, with every $a_i$ an integer, and rescale its tails by putting
\[
 T_N=2^{\,N+1}\sum_{i>N}\frac{a_i}{2^{\,i+1}}
    =\sum_{j\ge1}\frac{a_{N+j}}{2^{\,j}} .
\]
Two facts follow immediately.  First $T_{N+1}=2T_N-a_{N+1}$, so moving one
level along doubles the rescaled tail and subtracts a single coefficient.
Second $T_0=2\sum_{i\ge0}a_i2^{-(i+1)}-a_0$, so the sum is rational exactly
when $T_0$ is.  The whole of Section~\ref{sec:tail} therefore studies that
recurrence in isolation, assuming nothing about the coefficients beyond the
fact that they are integers; the prime-gap instance is resumed at the end of
the section.  The following definition is the object so obtained, stripped of
its origin.

\begin{definition}\label{def:rec}
Let $g:\N\to\Z$ and $T:\N\to\Q$.  Say $T$ satisfies the \emph{dyadic tail
recurrence} with \emph{digits} $g$ when
\[
 T_{N+1}=2T_N-g_{N+1}\qquad\text{for every }N .
\]
We call the sequence $(T_N)_{N\ge0}$ the \emph{orbit} of $T_0$ under $g$,
write $\sigma_h(N)=T_{N+h}-T_N$ for the \emph{shift} of length $h$ at $N$, and
call a rational number \emph{integral} when it is the image of an integer.
\end{definition}

\noindent These are the
\lword{Erdos251/PrimeGapDyadicTail.lean}{482}{DyadicTailRecurrence}{tail
recurrence}, the
\lword{Erdos251/PrimeGapDyadicTail.lean}{544}{tailShift}{shift}, and
\lword{Erdos251/PrimeGapDyadicTail.lean}{575}{RatIntegral}{integrality}.
The digits are arbitrary integers.  In the intended instance they are the prime
gaps and $T_N$ is the complete gap tail $\sum_{j\ge1}g_{N+j}2^{-j}$.  That
identification, and the recurrence over $\mathbb R$ with no rationality
hypothesis, already sit in \texttt{RealPrimeGapTail.lean}.  Every statement
below is a theorem about Definition~\ref{def:rec}.

One small orbit, referred to again below, is worth having in view.  Take every
digit $g_N=0$ and $T_0=1/12$.  Then $T_{N+1}=2T_N$, so the orbit is
$\tfrac1{12},\tfrac16,\tfrac13,\tfrac23,\tfrac43,\ldots$, and
$\sigma_h(N)=(2^{h}-1)2^{\,N}/12$.  The shift of length $2$ is not integral at
$N=0$, where $\sigma_2(0)=\tfrac13-\tfrac1{12}=\tfrac14$, and is integral at
$N=2$, where $\sigma_2(2)=\tfrac43-\tfrac13=1$; the shift of length $1$ is
$\sigma_1(N)=2^{\,N}/12$ and is never integral.  The change at $N=2$ is
accounted for by Theorem~\ref{res:collapse}, and the failure at $h=1$ by
Theorem~\ref{res:shiftiff}.

Iterating the recurrence $h$ times multiplies $T_N$ by $2^{h}$ and accumulates
an explicit integer, which we now name.  Define
$B_{0,N}=0$ and $B_{h+1,N}=2B_{h,N}+g_{N+h+1}$, so that
$B_{h,N}=g_{N+1}2^{\,h-1}+\cdots+g_{N+h}$: the
\lword{Erdos251/PrimeGapDyadicTail.lean}{647}{dyadicTailBlock}{tail block}.
Thus $B_{1,N}=g_{N+1}$, $B_{2,N}=2g_{N+1}+g_{N+2}$ and
$B_{3,N}=4g_{N+1}+2g_{N+2}+g_{N+3}$: the block puts the weights
$2^{\,h-1},\ldots,2^{0}$ on the $h$ digits following index $N$.

\begin{theorem}[block identity]\label{res:block}
Let $T:\N\to\Q$ satisfy the dyadic tail recurrence with integer digits $g$, and
let $B_{h,N}$ be as above.  For every $N$ and $h$,
\[
 T_{N+h}=2^{h}T_N-B_{h,N},
 \qquad\text{hence}\qquad
 \sigma_h(N)=(2^{h}-1)\,T_N-B_{h,N} .
\]
\end{theorem}

\begin{proof}
A routine induction on $h$; the step is one application of the recurrence together with
$2\cdot2^{h}=2^{h+1}$ and the recursion defining $B$.  The second identity is
the first minus $T_N$.
\end{proof}

\noindent Formalised as the
\lword{Erdos251/PrimeGapDyadicTail.lean}{653}{tail_iterate_eq_pow_mul_sub_block}{iterated
block identity} and the
\lword{Erdos251/PrimeGapDyadicTail.lean}{667}{tailShift_eq_scaled_sub_block}{scaled
shift identity}.  The shift also obeys the recurrence in its own right,
$\sigma_h(N+1)=2\sigma_h(N)-(g_{N+h+1}-g_{N+1})$, the
\lword{Erdos251/PrimeGapDyadicTail.lean}{563}{tailShift_succ}{shift step
identity}.

Since $B_{h,N}$ is an integer, Theorem~\ref{res:block} converts a question about
the shift into a question about the single scaled term $(2^{h}-1)T_N$.

\begin{theorem}[integral-shift criterion]\label{res:shiftiff}
Let $T:\N\to\Q$ satisfy the dyadic tail recurrence with integer digits $g$.
For every $N$ and $h$, the shift $\sigma_h(N)$ is integral if and only if
$(2^{h}-1)T_N$ is integral.
\end{theorem}

\begin{proof}
By Theorem~\ref{res:block} the two differ by the integer $B_{h,N}$, and
subtracting an integer does not change integrality.
\end{proof}

\noindent Formalised as the
\lword{Erdos251/PrimeGapDyadicTail.lean}{802}{tailShift_integral_iff_scaledTail}{integral-shift
criterion}, on the
\lword{Erdos251/PrimeGapDyadicTail.lean}{677}{ratIntegral_sub_int_iff}{integer-shift
invariance}.  This is an equivalence, not a one-way reduction: the block
carries no information about integrality, so the two conditions are the same
condition.  Informally, once $T_N$ is fixed no choice of the digits after index
$N$ can change whether the shift $\sigma_h(N)$ is an integer; that is decided
by the denominator of $T_N$ alone.

The denominator criterion is exact:
\[
  \sigma_h(N)\in\Z
  \quad\Longleftrightarrow\quad
  \operatorname{den}(T_N)\mid 2^h-1.
\tag{3.4}\label{eq:shift-denominator}
\]
This is the
\lword{Erdos251/PrimeGapDyadicTail.lean}{1279}{tailShift_integral_iff_den_dvd_mersenne}{denominator
classification}.  Euler's totient supplies one admissible shift length when
the denominator is odd; it is a witness, not the classification itself.

\begin{theorem}[a shift of totient length]\label{res:totient}
Let $T:\N\to\Q$ satisfy the dyadic tail recurrence with integer digits $g$, and
let $\ph$ be Euler's totient function.  If the reduced denominator $d$ of
$T_N$ is odd, then $\sigma_{\ph(d)}(N)$ is integral.
\end{theorem}

\begin{proof}
Since $d$ is odd, $2$ and $d$ are coprime, so Euler's congruence gives
$2^{\ph(d)}\equiv1\pmod d$, that is $d\mid 2^{\ph(d)}-1$.  Writing
$2^{\ph(d)}-1=dk$ and $T_N=u/d$ in lowest terms, $(2^{\ph(d)}-1)T_N=ku$ is an
integer, and Theorem~\ref{res:shiftiff} transfers this to the shift.
\end{proof}

\noindent Formalised as the
\lword{Erdos251/PrimeGapDyadicTail.lean}{877}{tailShift_integral_totient_of_odd_den}{totient
shift} and the
\lword{Erdos251/PrimeGapDyadicTail.lean}{695}{ratIntegral_totientMultiplier_of_odd_den}{Euler
multiplier}.  For example, if $\operatorname{den}T_N=3$ then $\ph(3)=2$ and
$(2^{2}-1)T_N=3T_N$ is an integer, so $\sigma_2(N)$ is integral, while
$(2^{1}-1)T_N=T_N$ is not; if $\operatorname{den}T_N=5$ then $\ph(5)=4$ and
$\sigma_4(N)$ is integral.

The hypothesis is a genuine restriction: the argument uses
coprimality of $2$ with the denominator, and the even part of a denominator is
exactly what the doubling in the recurrence acts on.  It cannot be dropped.
If $T_N=1/2$ then $2^{h}-1$ is odd for every $h\ge1$, so $(2^{h}-1)T_N$ is
never an integer and, by Theorem~\ref{res:shiftiff}, no shift at $N$ is
integral.

\begin{theorem}[propagation]\label{res:propagate}
Let $T:\N\to\Q$ satisfy the dyadic tail recurrence with integer digits $g$, and
fix $h$ and $N$.  If $\sigma_h(N)$ is integral, then $\sigma_h(N+k)$ is
integral for every $k\ge0$.
\end{theorem}

\begin{proof}
By the shift step identity, $\sigma_h(N+1)=2\sigma_h(N)-(g_{N+h+1}-g_{N+1})$ is
an integer combination of an integer and two digits; induct on $k$.
\end{proof}

\noindent Formalised as the
\lword{Erdos251/PrimeGapDyadicTail.lean}{888}{tailShift_integral_succ}{one-step
propagation} and the
\lword{Erdos251/PrimeGapDyadicTail.lean}{900}{tailShift_integral_add}{propagation
to every later index}.

The three preceding theorems combine as follows, and this is the statement the
rest of the note rests on.  The special case is immediate: if
$\operatorname{den}T_0$ is already odd, then Theorem~\ref{res:totient} at
$N=0$ makes the shift of length $\ph(\operatorname{den}T_0)$ integral and
Theorem~\ref{res:propagate} keeps it integral at every later index, so one may
take $N_0=0$.  In general a denominator carries a power of two as well, and the
key point is that the doubling in the recurrence annihilates
exactly the $2$-adic part of a denominator, and nothing else: after finitely
many steps the orbit therefore reaches a term with odd reduced denominator,
which is precisely the situation Theorem~\ref{res:totient} handles.  No
control of the digits is needed anywhere.

\begin{theorem}[exact denominator dynamics and eventual integrality]\label{res:collapse}
Let $T:\N\to\Q$ satisfy the dyadic tail recurrence with integer digits $g$
(Definition~\ref{def:rec}).  Then some fixed positive shift is integral from
some point onwards:
\[
 \exists h\ge1\ \exists N_0\ \forall N\ge N_0,\qquad
 \sigma_h(N)\in\Z .
\]
\end{theorem}

\begin{proof}
At every step the reduced denominator obeys the exact recurrence
\[
 \operatorname{den}(T_{N+1})
 =\frac{\operatorname{den}(T_N)}{\gcd(2,\operatorname{den}(T_N))}.
\]
Thus each even denominator loses exactly one factor of $2$, while an odd
denominator is unchanged.  After finitely many steps the denominator is odd.
Theorem~\ref{res:totient}, applied at that index $s$, supplies
the positive shift $h=\ph(\operatorname{den}T_s)$, and
Theorem~\ref{res:propagate} keeps that shift integral at every later index.
\end{proof}

\noindent The one-step formula is the
\lword{Erdos251/PrimeGapDyadicTail.lean}{1217}{tail_den_succ}{denominator recurrence},
with its
\lword{Erdos251/PrimeGapDyadicTail.lean}{1226}{tail_den_succ_eq_of_odd}{odd case}
and
\lword{Erdos251/PrimeGapDyadicTail.lean}{1236}{tail_den_succ_eq_half_of_even}{even case}.

In the orbit displayed after Definition~\ref{def:rec} the proof runs as
follows: $\operatorname{den}T_0=12=2^{2}\cdot3$, so $s=2$, the orbit reaches
$T_2=1/3$ with odd denominator, and $h=\ph(3)=2$.  That is exactly the shift
length seen to be integral there from index $2$ onwards, and no shorter one
works.

Three features of the argument are used later.  The proof may take as its
number of preparatory steps the $2$-adic valuation of
$\operatorname{den}T_0$.  The resulting shift length $h$ is the totient of the
odd denominator reached from a hypothetical rational initial value, and is
not known in advance for the prime-gap orbit.  This is why this argument
requires Problem~\ref{prob:escape} for every $h$, rather than for one
preassigned shift length.  Finally, the digits enter only through the integer
$B_{s,0}$, so the conclusion holds for an arbitrary integer digit sequence.

\noindent In the current formal source, denominator factorisation and
cancellation are packaged directly in the
\lword{Erdos251/PrimeGapDyadicTail.lean}{830}{rationalPrimeGapTailShift_eventuallyIntegral_of_den_eq}{fixed-denominator
theorem}, and the quantified conclusion is the
\lword{Erdos251/PrimeGapDyadicTail.lean}{853}{rationalPrimeGapTailShift_eventuallyIntegral}{eventual
integral-shift theorem} for the actual prime-gap tail state.

The useful contrapositive is stated for a real recurrence.  Call its shifts
\emph{cofinally non-integral} when, for every fixed $h\ge1$ and every
threshold $N_0$, some $N\ge N_0$ has $\sigma_h(N)\notin\Z$.  This is precisely
the negation of the conclusion of Theorem~\ref{res:collapse}: no shift length
whatever becomes integral and stays integral.

\begin{theorem}[exact rationality classification]\label{res:escape-irrational}
Let $T:\N\to\R$ satisfy $T_{N+1}=2T_N-g_{N+1}$ with integer $g$.  If its
shifts are defined by $\sigma_h(N)=T_{N+h}-T_N$, then the following are
equivalent:
\begin{enumerate}[label=\textup{(\roman*)}]
\item $T_0$ is rational;
\item $\sigma_h(N)$ is integral for some $h\ge1$ and some $N$;
\item for some fixed $h\ge1$, $\sigma_h(N)$ is integral at every sufficiently
  large $N$.
\end{enumerate}
Consequently $T_0$ is irrational if and only if every positive-length shift
is non-integral at every index, equivalently if and only if the shifts are
cofinally non-integral.
\end{theorem}

\begin{proof}
For \textup{(i)}$\Rightarrow$\textup{(iii)}, choose $q\in\Q$ whose real cast
is $T_0$.  The real block identity identifies the whole orbit with the cast of
the rational recurrence starting at $q$; Theorem~\ref{res:collapse} applied to
that rational orbit then gives \textup{(iii)}.  The implication
\textup{(iii)}$\Rightarrow$\textup{(ii)} is immediate.  For
\textup{(ii)}$\Rightarrow$\textup{(i)}, the real block identity gives
\[
 \sigma_h(N)=(2^h-1)T_N-B_{h,N}.
\]
Here $B_{h,N}$ and $\sigma_h(N)$ are integers and $2^h-1\ne0$, so $T_N$ is
rational.  Iterating the recurrence backwards through the block identity then
makes $T_0$ rational.  Negating the pointwise and eventual forms gives the two
irrationality formulations.
\end{proof}

\noindent Lean checks the rational actual-tail state as
\lword{Erdos251/PrimeGapDyadicTail.lean}{489}{rationalPrimeGapTailState}{rational
prime-gap tail state}, its
\lword{Erdos251/PrimeGapDyadicTail.lean}{527}{rationalPrimeGapTailState_recurrence}{recurrence},
and the bridge from a hypothetical rational value to that state as
\lword{Erdos251/PrimeGapDyadicTail.lean}{510}{exists_rationalPrimeGapTailState_representation_of_not_irrational}{the
rational-tail representation}.  The exact real classifiers are
\lword{Erdos251/PrimeGapDyadicTail.lean}{1479}{not_irrational_initial_iff_exists_integral_positive_tailShift}{one
integral positive shift},
\lword{Erdos251/PrimeGapDyadicTail.lean}{1525}{not_irrational_initial_iff_exists_eventually_integral_positive_tailShift}{eventual
integrality},
\lword{Erdos251/PrimeGapDyadicTail.lean}{1551}{irrational_initial_iff_all_positive_tailShifts_nonintegral}{pointwise
non-integrality}, and
\lword{Erdos251/PrimeGapDyadicTail.lean}{1572}{irrational_initial_iff_cofinalNonintegralTailShifts}{cofinal
non-integrality}.

\paragraph{The actual prime-gap orbit.}
The concrete tails
\[
 T_N=\sum_{j\ge1}\frac{g_{N+j}}{2^j}
\]
are already identified in Lean, with no rationality hypothesis
(\lword{Erdos251/RealPrimeGapTail.lean}{31}{realPrimeGapTail_eq_tsum_shifted_gaps}{shifted-gap tail};
\lword{Erdos251/RealPrimeGapTail.lean}{48}{realPrimeGapTail_recurrence}{actual recurrence};
\lword{Erdos251/RealPrimeGapTail.lean}{73}{irrational_primeSeries_iff_realPrimeGapTail_escape}{escape equivalence}).
The Lean-checked polynomial prime bound gives convergence, and
$T_{N+1}=2T_N-g_{N+1}$ is the same identification.  If
$G=\sum_{n\ge0}g_n/2^{n+1}$, then $T_0=2G-1$.  Together with
Theorem~\ref{res:infinite}, this shows that $\Pi$, $G$, and
$T_0$ have the same rationality status.  Thus
Theorem~\ref{res:escape-irrational} identifies Problem~\#251 exactly with
cofinal shift escape for $T$.  The older rational-candidate representation
remains in \texttt{PrimeGapDyadicTail.lean}.  What is not
proved is the cofinal non-integrality needed by
Theorem~\ref{res:escape-irrational}.

\subsection{Two adjacent small shifts cannot both be integral}\label{sec:local-certificate}

The exact classifier reduces irrationality to cofinal non-integrality, a
condition of infinite precision imposed on a complete tail.  The key
point is that inside the open interval $(-1,1)$ integrality is equality with
zero, so on that range the one-step shift recurrence
\[
 \sigma_h(N+1)=2\sigma_h(N)-\bigl(g_{N+h+1}-g_{N+1}\bigr)
\]
turns simultaneous integrality of two adjacent shifts into a single comparison
of digits.  What this buys is a \emph{certificate}: a condition attached to a
single pair of adjacent indices, whose verification already contradicts
integrality at that pair.  The point of arranging the argument this
way is that the distance of a shift from the integers never has to be
estimated; it suffices to know that both shifts lie in $(-1,1)$ and that two
digits differ.

\begin{theorem}[adjacent small-shift obstruction]\label{res:smallpair}
Let $T:\N\to\Q$ satisfy the dyadic tail recurrence with integer digits $g$.
Fix $h$ and $N$.  If
\[
 -1<\sigma_h(N)<1,\qquad -1<\sigma_h(N+1)<1,
 \qquad g_{N+h+1}\ne g_{N+1},
\]
then $\sigma_h(N)$ and $\sigma_h(N+1)$ cannot both be integral.
Consequently, if such a pair occurs beyond every threshold, the $h$-shift is
not eventually integral.
\end{theorem}

\begin{proof}
An integral rational strictly between $-1$ and $1$ is zero.  If both shifts
were integral, both would therefore vanish, and substitution in the displayed
shift step identity would give $g_{N+h+1}=g_{N+1}$, a contradiction.  The
cofinal statement chooses one such adjacent pair after the alleged onset of
integrality.
\end{proof}

\noindent The finite contradiction is the
\lword{Erdos251/PrimeGapDyadicTail.lean}{979}{tailShift_not_both_integral_of_small_pair_of_digit_ne}{adjacent
small-shift obstruction}; its quantified form is the
\lword{Erdos251/PrimeGapDyadicTail.lean}{1006}{not_eventuallyIntegralTailShift_of_cofinal_small_mismatch}{cofinal
small-mismatch theorem}; and the actual-prime-gap specialisation is the
\lword{Erdos251/PrimeGapDyadicTail.lean}{1112}{primeGapTailShift_not_eventuallyIntegral_of_cofinal_small_mismatch}{prime-gap
specialisation}.  Theorem~\ref{res:smallpair} is conditional on its two tail
inequalities; no theorem asserting that such pairs occur is claimed here.

All three are stated for a rational orbit, while the prime-gap tail
$T$ of Section~\ref{sec:tail} is real, so the rational statement does
not on its own discharge the instances that arise downstream.  The same
argument gives the real form directly.

\begin{corollary}[adjacent small-shift obstruction, real form]\label{res:smallpair-real}
Let $T:\N\to\R$ satisfy $T_{N+1}=2T_N-g_{N+1}$ with integer digits $g$, and
write $\sigma_h(N)=T_{N+h}-T_N$.  Fix $h$ and $N$.  If
\[
 -1<\sigma_h(N)<1,\qquad -1<\sigma_h(N+1)<1,
 \qquad g_{N+h+1}\ne g_{N+1},
\]
then $\sigma_h(N)$ and $\sigma_h(N+1)$ do not both lie in $\Z$.
Consequently, if such a pair occurs beyond every threshold, the $h$-shift is
not eventually integral.
\end{corollary}

\begin{proof}
The real recurrence gives the same shift step identity
$\sigma_h(N+1)=2\sigma_h(N)-\bigl(g_{N+h+1}-g_{N+1}\bigr)$, and an integer
strictly between $-1$ and $1$ is zero.  If both shifts lay in $\Z$, both would
therefore vanish, and substitution in that identity would give
$g_{N+h+1}=g_{N+1}$, a contradiction.  The cofinal statement chooses one such
adjacent pair after the alleged onset of integrality.
\end{proof}

\noindent Corollary~\ref{res:smallpair-real} and the resulting implication
for the actual prime series are proved in the live module
\texttt{RealPrimeGapTail.lean}.  These live declarations can postdate the
historical public snapshot used by the source links.

The third hypothesis, on its own, is available.  For the actual gaps tabulated
in Section~\ref{sec:problem} it reads $g_4=2\ne4=g_3$ at $h=1$, $N=2$; and
Proposition~\ref{res:gap-nonperiodic} below says precisely that for each fixed
$h\ge1$ the inequality $g_{N+h+1}\ne g_{N+1}$ holds for arbitrarily large $N$,
since its failure from some index onwards is eventual periodicity with period
$h$.  What is missing is their joint occurrence with the digit mismatch at the
same indices.  Each inequality constrains a complete infinite tail.

\begin{proposition}[prime gaps do not become periodic]\label{res:gap-nonperiodic}
For every positive $h$, the actual consecutive-prime-gap sequence is not
eventually periodic with period $h$.
\end{proposition}

\begin{proof}
The gaps are unbounded, by the standard construction: the interval from $n!+2$
to $n!+n$ contains no prime.  Far stronger lower bounds for large gaps are
known~\cite[Theorem~1, p.~66]{fgkmt2018}, but unboundedness is all that is
needed.  An
eventually periodic natural-valued sequence has finite range after its
preperiod, while its finite initial segment is bounded as well; hence it is
bounded, a contradiction.
\end{proof}

\noindent Lean checks the factorial argument as
\lword{Erdos251/PrimeGapDyadicTail.lean}{57}{exists_primeGap0_gt}{unboundedness
of the actual gaps} and the conclusion as
\lword{Erdos251/PrimeGapDyadicTail.lean}{1023}{primeGap0_not_eventually_periodic}{non-eventual
periodicity}.  Combining nonperiodicity with eventual strict smallness of one
positive shift would also exclude eventual integrality, but eventual
smallness at every sufficiently large index is stronger than the cofinal
adjacent-pair hypothesis in Theorem~\ref{res:smallpair} and is not asserted
here.

\section{Nonperiodic coefficients with a rational sum}\label{sec:carry}

Proposition~\ref{res:gap-nonperiodic} shows that the prime gaps are not
eventually periodic.  One might therefore hope that rationality of a dyadic
series forces its integer coefficients to be eventually periodic, and play the
two against each other.  The hoped-for implication is false, and a single
explicit sequence refutes it.

The construction runs the emission of coefficients backwards.  Let
$K:\N\to\Q$ be arbitrary, read $K_n$ as the value carried into level $n$, and
put $\kappa_n=2K_n-K_{n+1}$: the
\lword{Erdos251/PrimeGapDyadicTail.lean}{1129}{carryCoeff}{carry coefficient}.
That definition is exactly the statement that
\[
 \frac{K_n}{2^{\,n}}
 =\frac{\kappa_n}{2^{\,n+1}}+\frac{K_{n+1}}{2^{\,n+1}} ,
\]
so at each level the carried value splits into one emitted coefficient and a
new carry.  Nothing at all is assumed about $K$: not integrality, not
positivity, not any bound.  That is the sense in which the carry is free, and
it is what the counterexample below exploits.

\begin{proposition}[exact telescoping]\label{res:telescope}
Let $K:\N\to\Q$ be arbitrary and $\kappa_n=2K_n-K_{n+1}$.  For every $n\ge0$,
\[
 \sum_{i=0}^{n-1}\frac{\kappa_i}{2^{\,i+1}}=K_0-\frac{K_n}{2^{\,n}} .
\]
\end{proposition}

\begin{proof}
A routine induction on $n$; the added term is
$(2K_n-K_{n+1})/2^{\,n+1}=K_n/2^{\,n}-K_{n+1}/2^{\,n+1}$.
\end{proof}

\noindent Formalised as the
\lword{Erdos251/PrimeGapDyadicTail.lean}{1137}{carryPartialSum_eq}{carry
telescoping identity}, on the
\lword{Erdos251/PrimeGapDyadicTail.lean}{1133}{carryPartialSum}{carry partial
sum}.  When $K$ takes natural-number values and $K_{n+1}\le2K_n$ for every $n$,
the coefficients $\kappa_n$ are natural numbers as well, and the
natural-number and rational readings of the definition agree under the cast,
the
\lword{Erdos251/PrimeGapDyadicTail.lean}{1180}{natCarryCoeff_cast}{natural-carry
cast}.

\paragraph{Consequence for the coefficients.}
If $K_n2^{-n}\to0$, the emitted partial sums converge to $K_0$.  A
parity-compatible positive example is
\[
 K_0=\frac52,\qquad K_n=2n+2\ (n\ge1),\qquad
 \kappa_0=1,\quad\kappa_n=2n\ (n\ge1),
\]
for which
\[
 \sum_{i=0}^{n-1}\frac{\kappa_i}{2^{\,i+1}}
   =\frac52-\frac{K_n}{2^n}\longrightarrow\frac52.
\]
Thus the emitted coefficients $1,2,4,6,8,10,\ldots$ are positive, even after
the first term, unbounded and not eventually periodic, although their dyadic
sum is rational.  Rationality alone therefore cannot imply eventual
periodicity even for a positive, parity-correct integer coefficient sequence.
The example does not claim that these coefficients are prime gaps; it isolates
the additional arithmetic information any successful argument must use.

\section{Complements and further questions}\label{sec:open}

Problem~\#251 is open.  The public Lean source now proves unconditional
convergence of both dyadic series, their exact infinite summation-by-parts
identity, and the real-to-rational scaled-tail representation
(\lword{Erdos251/PrimeGapDyadicTail.lean}{427}{tsum_primeDyadicTerm_eq_two_add_primeGap_unconditional}{identity},
\lword{Erdos251/PrimeGapDyadicTail.lean}{510}{exists_rationalPrimeGapTailState_representation_of_not_irrational}{tail representation}).
It also proves that every rational candidate supplies one positive fixed shift
which is integral at every sufficiently late tail index
(\lword{Erdos251/PrimeGapDyadicTail.lean}{853}{rationalPrimeGapTailShift_eventuallyIntegral}{checked}).
The denominator construction selects that same shift so that prime-gap
nonperiodicity also rules out its eventual confinement to the open unit
interval
(\lword{Erdos251/PrimeGapDyadicTail.lean}{1097}{rationalPrimeGapTail_has_positive_shift_not_eventually_small}{not eventually small}).
This conclusion is compatible with rationality: it says that the rationally
forced shift cannot support an eventual-smallness contradiction.  It supplies
neither eventual smallness nor cofinally many adjacent small mismatches for
that shift.
The remaining issue is recurrence or anti-concentration of the \emph{actual}
prime-gap shifts.  Throughout, $g_i=p_{i+1}-p_i$ are the zero-based prime gaps
of Section~\ref{sec:problem}.  Write $\mathsf E(h)$ for the cofinal escape
statement in~\eqref{eq:shift-escape} at shift $h$.

\begin{problem}[divisor-hitting shift escape]\label{prob:divisor-hit}
For every $r\ge1$, does some positive multiple of $r$ escape cofinally?
\[
  \forall r\ge1\ \exists m\ge1:\qquad \mathsf E(mr).
\]
\end{problem}

\noindent Divisor-hitting escape changes the organisation of the
quantifiers.  Within an integer-digit recurrence it is equivalent to
irrationality, just as universal shift escape is: a rational state has a
positive eventual order, and every positive multiple of that order is an
integral shift.  A quantitatively stated sufficient estimate can still be
easier to attack, but that requires a separately proved analytic bound.
Forward propagation is Lean-checked.

\begin{problem}[universal prime-gap shift escape]\label{prob:escape}
For every $h\ge1$ and every $N_0$, prove that some $N\ge N_0$ satisfies
\[
 \sum_{j\ge1}
 \frac{g_{N+h+j}-g_{N+j}}{2^j}\notin\Z .
\tag{5.1}\label{eq:shift-escape}
\]
\end{problem}

The series in~\eqref{eq:shift-escape} is exactly $T_{N+h}-T_N$.  Theorem~\ref{res:escape-irrational} makes
Problem~\ref{prob:escape} equivalent to irrationality of $\Pi$.  Within the
integer-digit recurrence, divisor-hitting and universal escape are equivalent:
a hypothetical rational value chooses a shift length from the odd part of its
reduced denominator, and the cocycle then makes every positive multiple
eventually integral.  The two formulations organise the quantifiers
differently; they do not give a strict logical weakening.

\begin{problem}[cofinal adjacent small mismatch]\label{prob:smallpair}
For every fixed $h\ge1$ and every $N_0$, prove that some $N\ge N_0$ satisfies
\[
\begin{split}
\left|\sum_{j\ge1}
 \frac{g_{N+h+j}-g_{N+j}}{2^j}\right|&<1,\\
\left|\sum_{j\ge1}
 \frac{g_{N+h+1+j}-g_{N+1+j}}{2^j}\right|&<1,\\
g_{N+h+1}&\ne g_{N+1}.
\end{split}
\tag{5.2}\label{eq:smallpair}
\]
\end{problem}

The two sums are adjacent $h$-shifts.  Corollary~\ref{res:smallpair-real} turns
each instance of~\eqref{eq:smallpair} into a finite contradiction to
simultaneous integrality; Theorem~\ref{res:smallpair} is the Lean-checked
rational case.  A cofinal family of such pairs therefore proves
Problem~\ref{prob:escape}; Theorem~\ref{res:escape-irrational} then gives irrationality
of the gap series, and Corollary~\ref{res:irr-equivalence} transfers it to
$\Pi$.  The endpoint takes every positive $h$: a proof only at $h=1$
eliminates one eventual-integrality possibility and does not discharge the
classifier.  Digit mismatches and small shifts must occur at the same indices;
separate cofinal occurrence statements do not establish that intersection.
This formulation deliberately asks only for sporadic adjacent pairs;
the stronger assertion that a fixed shift is eventually always smaller than
one is unnecessary.  The conjunction may have density zero: that excludes a
positive-proportion lower bound, but an averaging argument may still produce
an unbounded sparse count.

Several natural prime-distribution inputs are insufficient.  Isolated
small gaps, isolated large gaps, average gap estimates, and the occurrence of
any one fixed finite pattern do not suffice: both inequalities in
\eqref{eq:smallpair} contain the complete infinite continuation.  Nor
does parity: after the first gap all $g_n$ are even.  Unboundedness and
non-eventual-periodicity of the actual gaps, though now checked, do not imply
the required small-tail recurrence.  In particular, Section~\ref{sec:carry}
shows that no argument can deduce eventual periodicity from rationality alone.

\paragraph{A finite truncation criterion.}
Both problems above are stated in terms of infinite tails, but a dominated
truncation reduces the first of them to a finite quantity.  For $L\ge1$ put
\[
 S_{h,N,L}=\sum_{j=1}^{L}
 \frac{g_{N+h+j}-g_{N+j}}{2^j},
\]
the truncation of the series in~\eqref{eq:shift-escape} after $L$ terms.  This
is a finite sum of gap differences and can be computed; the point of the
following proposition is to say how far from an integer it must be before the
discarded tail is irrelevant.

Equivalently, introduce the integral dyadic block
\[
 D_{h,N,L}=\sum_{j=1}^{L}2^{L-j}
   (g_{N+h+j}-g_{N+j}),\qquad S_{h,N,L}=\frac{D_{h,N,L}}{2^L}.
\]
Then
\[
 \operatorname{dist}(S_{h,N,L},\Z)
 =2^{-L}\min\{D_{h,N,L}\bmod2^L,
 2^L-(D_{h,N,L}\bmod2^L)\}.
\]
Here $D_{h,N,L}\bmod2^L$ denotes the least nonnegative residue, including
when $D_{h,N,L}<0$.
Thus the finite criterion below is an exact modular small-arc problem: the
residue of $D_{h,N,L}$ must avoid the two arcs of radius
$2^LR_{h,N,L}(M)$ around $0$ modulo $2^L$.  A one-block certificate would be a
prime-gap theorem producing such an avoided arc on a logarithmic block.

\begin{proposition}[finite truncation]\label{res:truncation}
Let $M(n)\ge g_n$ for every $n$, assume that the series below converges, and
put
\[
 R_{h,N,L}(M)=
 \sum_{j>L}\frac{M(N+h+j)+M(N+j)}{2^j}.
\]
Suppose that for every fixed $h\ge1$ and every $N_0$ there exist $N\ge N_0$ and
$L\ge1$ with
\[
 \operatorname{dist}(S_{h,N,L},\Z)>R_{h,N,L}(M).
\tag{5.3}\label{eq:truncation}
\]
Then Problem~\ref{prob:escape} holds.
\end{proposition}

\begin{proof}
The part of $\sum_{j\ge1}(g_{N+h+j}-g_{N+j})2^{-j}$ omitted from $S_{h,N,L}$
has absolute value at most $R_{h,N,L}(M)$, so under~\eqref{eq:truncation} the
full sum lies at positive distance from every integer.
\end{proof}

\noindent The last distance-to-integers inference is the elementary paper
argument in the displayed proof: an error at most $R$ cannot reach an integer
when the approximation is farther than $R$ from every integer.  It is not
currently a named Lean declaration.  For instance, if
$S_{h,N,L}=3/8$ and $R_{h,N,L}(M)=1/32$, the full sum lies within $1/32$ of
$3/8$ and so at distance at least $11/32$ from every integer, which settles
\eqref{eq:shift-escape} at that $N$.  The prime-gap tail
bound, the convergence used above, and the existence of blocks satisfying
\eqref{eq:truncation} are not consequences of that inference.  The
proposition's full sum is moreover real-valued, so the reverse-triangle step
used here is a paper proof and not an instance of any rational-valued formal
declaration.

\noindent Taking the classical bound $M(n)\ll n\log n$, a choice
$L=\lceil A\log_2(N+h+2)\rceil$ with any fixed $A>1$ makes the right side a
negative power of $N$ up to logarithms.  Thus~\eqref{eq:truncation} asks for a finite
dyadic anti-concentration estimate on a logarithmic-length block, not control
of an infinite tail and not eventual periodicity of the full gap sequence.
For Problem~\ref{prob:smallpair}, the same truncation must certify two adjacent
full-tail values inside the open unit interval, together with the displayed
gap mismatch.  A finite prefix is useful only when its omitted tail is
rigorously dominated.

What~\eqref{eq:truncation} requires is joint control of the finite block of
weighted differences $(g_{N+h+1}-g_{N+1},\ldots,g_{N+h+L}-g_{N+L})$ modulo
powers of two, along a block of logarithmic length.  We do not know how to obtain such control and
we make no progress on it here.  The strongest results on prime gaps address a
different shape of question.  Zhang's bounded-gap theorem
\cite[Theorem~1, p.~1122]{zhang2014} produces infinitely many bounded
consecutive-prime gaps.
Maynard proves substantially more than an individual-gap statement:
\cite[Theorem~1.1, p.~384]{maynard2015} bounds
$\liminf_n(p_{n+m}-p_n)$ for every fixed $m$, and hence gives bounded clusters
of every fixed size; \cite[Theorem~1.3, p.~385]{maynard2015} gives the explicit
unconditional bound $\liminf_n(p_{n+1}-p_n)\le 600$.  The large-gap theorem of
Ford, Green, Konyagin, Maynard and Tao
\cite[Theorem~1, p.~66]{fgkmt2018} bounds the largest single consecutive-prime
gap below $X$.  None of these results supplies the joint dyadic distribution
of a logarithmic block of consecutive gap differences required by
\eqref{eq:truncation}.  The same introduction notes a separate sequel on
chains of large gaps; the cited theorem itself supplies no residue-sensitive
block estimate of the kind needed here.

\paragraph{Established bridges and the remaining estimate.}
The live source contains convergence, the prime-to-gap identity, the actual
real-tail recurrence, the rationality classifier, and the conditional
small-mismatch implication.  The unproved assertion is the supply of the
joint events in~\eqref{eq:smallpair}, not the construction of the tail
itself.  The ordinary mean-tail and dispersion criteria recorded in the
companion reasoning record have unproved lower-bound hypotheses as well.
A finite continued-fraction enclosure separately gives
$q\ge2^{39997}>10^{12040}$ for any rational representation of $\Pi$;
this is a finite denominator exclusion, not an irrationality measure.

\pdfbookmark[1]{Statements and declarations}{statements-declarations}
\subsection*{Statements and declarations}

\paragraph{Artefact and data availability.} The
\href{\sourceurl}{source links} identify a fixed public snapshot.
More recent live declarations are identified explicitly where they are used;
a local declaration locator does not establish its presence at a historical
public pin.  The repository's versioned source and verification records,
rather than this manuscript's navigation links, identify the checked object.

Lean checks each proof term against the fixed library
version, and the sources linked here contain no proof placeholders and no
project-defined axioms; Lean does not authorise the exposition, the citation
choices, or the interpretation, for which the author remains responsible.

\paragraph{Funding and competing interests.} This work received no external
funding.  The author declares no competing interests.

\paragraph{Acknowledgements.} The problem numbering and status follow the
Erd\H{o}s Problems catalogue maintained by Thomas Bloom~\cite{erdosproblems}.

\appendix
\section{Guide to the formal sources}\label{app:index}

Each linked phrase opens its Lean declaration at the pinned source
revision~\commitshort.  The note's declarations live in the \#251 modules,
including the actual real-tail bridge.  A local locator is not evidence that
the historical public pin contains the file.  The summation-by-parts
declarations are prime-specific; most of Section~\ref{sec:tail} is stated for
arbitrary integer digits and arbitrary rational or real orbits, while
Section~\ref{sec:local-certificate} records the actual-gap specialisation.

% The exhaustive source manifest is retained in the authored TeX for automated
% validation, and kept outside the printed note.
\iffalse

\begin{tabular}{@{}ll@{}}
informal statement & linked Lean source\\
zero-based primes & \lrefx{Erdos251/PrimeGapDyadicTail.lean}{40}{prime0}\\
zero-based gaps & \lrefx{Erdos251/PrimeGapDyadicTail.lean}{44}{primeGap0}\\
first gap value & \lrefx{Erdos251/PrimeGapDyadicTail.lean}{47}{primeGap0_zero}\\
second gap value & \lrefx{Erdos251/PrimeGapDyadicTail.lean}{51}{primeGap0_one}\\
unbounded actual prime gaps & \lrefx{Erdos251/PrimeGapDyadicTail.lean}{57}{exists_primeGap0_gt}\\
dyadic partial sum & \lrefx{Erdos251/PrimeGapDyadicTail.lean}{101}{dyadicPartialSumQ}\\
displayed normalisation & \lrefx{Erdos251/PrimeGapDyadicTail.lean}{106}{prime0DisplayedPartialSumQ}\\
factor-of-two identity & \lrefx{Erdos251/PrimeGapDyadicTail.lean}{111}{prime0DisplayedPartialSumQ_eq_two_mul}\\
dyadic difference sum & \lrefx{Erdos251/PrimeGapDyadicTail.lean}{121}{dyadicDifferencePartialSumQ}\\
partial-sum step & \lrefx{Erdos251/PrimeGapDyadicTail.lean}{124}{dyadicPartialSumQ_succ}\\
difference-sum step & \lrefx{Erdos251/PrimeGapDyadicTail.lean}{129}{dyadicDifferencePartialSumQ_succ}\\
summation by parts & \lrefx{Erdos251/PrimeGapDyadicTail.lean}{138}{dyadicPartialSumQ_eq_start_add_differences}\\
monotonicity step & \lrefx{Erdos251/PrimeGapDyadicTail.lean}{152}{prime0_mono_step}\\
gap cast identity & \lrefx{Erdos251/PrimeGapDyadicTail.lean}{156}{primeGap0_cast}\\
gap partial sum & \lrefx{Erdos251/PrimeGapDyadicTail.lean}{161}{primeGapPartialSumQ}\\
prime-gap summation by parts & \lrefx{Erdos251/PrimeGapDyadicTail.lean}{172}{prime0_dyadic_summation_by_parts}\\
prime dyadic term & \lrefx{Erdos251/PrimeGapDyadicTail.lean}{183}{primeDyadicTerm}\\
displayed prime dyadic term & \lrefx{Erdos251/PrimeGapDyadicTail.lean}{188}{primeDisplayedDyadicTerm}\\
prime-gap dyadic term & \lrefx{Erdos251/PrimeGapDyadicTail.lean}{192}{primeGapDyadicTerm}\\
infinite factor-of-two identity & \lrefx{Erdos251/PrimeGapDyadicTail.lean}{196}{primeDisplayedDyadicTerm_eq_two_mul}\\
gap discrete derivative & \lrefx{Erdos251/PrimeGapDyadicTail.lean}{202}{primeGapDyadicTerm_eq}\\
gap summability transfer & \lrefx{Erdos251/PrimeGapDyadicTail.lean}{385}{summable_primeGapDyadicTerm_of_summable_primeDyadicTerm}\\
infinite gap identity & \lrefx{Erdos251/PrimeGapDyadicTail.lean}{404}{tsum_primeDyadicTerm_eq_two_add_primeGap}\\
normalised irrationality equivalence & \lrefx{Erdos251/PrimeGapDyadicTail.lean}{435}{irrational_tsum_primeDyadicTerm_iff_primeGap}\\
displayed infinite identity & \lrefx{Erdos251/PrimeGapDyadicTail.lean}{444}{tsum_primeDisplayedDyadicTerm_eq_four_add_two_primeGap}\\
displayed irrationality equivalence & \lrefx{Erdos251/PrimeGapDyadicTail.lean}{459}{irrational_tsum_primeDisplayedDyadicTerm_iff_primeGap}\\
tail recurrence & \lrefx{Erdos251/PrimeGapDyadicTail.lean}{482}{DyadicTailRecurrence}\\
shift & \lrefx{Erdos251/PrimeGapDyadicTail.lean}{544}{tailShift}\\
shift step identity & \lrefx{Erdos251/PrimeGapDyadicTail.lean}{563}{tailShift_succ}\\
integrality & \lrefx{Erdos251/PrimeGapDyadicTail.lean}{575}{RatIntegral}\\
tail block & \lrefx{Erdos251/PrimeGapDyadicTail.lean}{647}{dyadicTailBlock}\\
iterated block identity & \lrefx{Erdos251/PrimeGapDyadicTail.lean}{653}{tail_iterate_eq_pow_mul_sub_block}\\
scaled shift identity & \lrefx{Erdos251/PrimeGapDyadicTail.lean}{667}{tailShift_eq_scaled_sub_block}\\
integer-shift invariance & \lrefx{Erdos251/PrimeGapDyadicTail.lean}{677}{ratIntegral_sub_int_iff}\\
Euler multiplier & \lrefx{Erdos251/PrimeGapDyadicTail.lean}{695}{ratIntegral_totientMultiplier_of_odd_den}\\
integral-shift criterion & \lrefx{Erdos251/PrimeGapDyadicTail.lean}{802}{tailShift_integral_iff_scaledTail}\\
totient shift & \lrefx{Erdos251/PrimeGapDyadicTail.lean}{877}{tailShift_integral_totient_of_odd_den}\\
one-step propagation & \lrefx{Erdos251/PrimeGapDyadicTail.lean}{888}{tailShift_integral_succ}\\
propagation to every later index & \lrefx{Erdos251/PrimeGapDyadicTail.lean}{900}{tailShift_integral_add}\\
rational actual-tail state & \lrefx{Erdos251/PrimeGapDyadicTail.lean}{489}{rationalPrimeGapTailState}\\
rational actual-tail recurrence & \lrefx{Erdos251/PrimeGapDyadicTail.lean}{527}{rationalPrimeGapTailState_recurrence}\\
hypothetical-rational tail representation & \lrefx{Erdos251/PrimeGapDyadicTail.lean}{510}{exists_rationalPrimeGapTailState_representation_of_not_irrational}\\
power-of-two/odd denominator scaling & \lrefx{Erdos251/PrimeGapDyadicTail.lean}{723}{ratIntegral_scaled_of_den_eq_pow_two_mul_odd}\\
multiplier at the two-adic exponent & \lrefx{Erdos251/PrimeGapDyadicTail.lean}{756}{ratIntegral_multiplier_tail_at_twoExponent}\\
fixed-denominator eventual shift & \lrefx{Erdos251/PrimeGapDyadicTail.lean}{830}{rationalPrimeGapTailShift_eventuallyIntegral_of_den_eq}\\
eventual integral shift & \lrefx{Erdos251/PrimeGapDyadicTail.lean}{853}{rationalPrimeGapTailShift_eventuallyIntegral}\\
adjacent small-shift obstruction & \lrefx{Erdos251/PrimeGapDyadicTail.lean}{979}{tailShift_not_both_integral_of_small_pair_of_digit_ne}\\
cofinal small-mismatch theorem & \lrefx{Erdos251/PrimeGapDyadicTail.lean}{1006}{not_eventuallyIntegralTailShift_of_cofinal_small_mismatch}\\
prime gaps not eventually periodic & \lrefx{Erdos251/PrimeGapDyadicTail.lean}{1023}{primeGap0_not_eventually_periodic}\\
actual-prime-gap specialisation & \lrefx{Erdos251/PrimeGapDyadicTail.lean}{1112}{primeGapTailShift_not_eventuallyIntegral_of_cofinal_small_mismatch}\\
carry coefficient & \lrefx{Erdos251/PrimeGapDyadicTail.lean}{1129}{carryCoeff}\\
carry partial sum & \lrefx{Erdos251/PrimeGapDyadicTail.lean}{1133}{carryPartialSum}\\
carry telescoping identity & \lrefx{Erdos251/PrimeGapDyadicTail.lean}{1137}{carryPartialSum_eq}\\
natural carry coefficient & \lrefx{Erdos251/PrimeGapDyadicTail.lean}{1177}{natCarryCoeff}\\
natural-carry cast & \lrefx{Erdos251/PrimeGapDyadicTail.lean}{1180}{natCarryCoeff_cast}\\
exact denominator recurrence & \lrefx{Erdos251/PrimeGapDyadicTail.lean}{1217}{tail_den_succ}\\
odd-denominator persistence & \lrefx{Erdos251/PrimeGapDyadicTail.lean}{1226}{tail_den_succ_eq_of_odd}\\
even-denominator halving & \lrefx{Erdos251/PrimeGapDyadicTail.lean}{1236}{tail_den_succ_eq_half_of_even}\\
integral shifts by denominator & \lrefx{Erdos251/PrimeGapDyadicTail.lean}{1279}{tailShift_integral_iff_den_dvd_mersenne}\\
one-shift rationality classifier & \lrefx{Erdos251/PrimeGapDyadicTail.lean}{1479}{not_irrational_initial_iff_exists_integral_positive_tailShift}\\
eventual-shift rationality classifier & \lrefx{Erdos251/PrimeGapDyadicTail.lean}{1525}{not_irrational_initial_iff_exists_eventually_integral_positive_tailShift}\\
pointwise irrationality classifier & \lrefx{Erdos251/PrimeGapDyadicTail.lean}{1551}{irrational_initial_iff_all_positive_tailShifts_nonintegral}\\
cofinal irrationality classifier & \lrefx{Erdos251/PrimeGapDyadicTail.lean}{1572}{irrational_initial_iff_cofinalNonintegralTailShifts}\\
\end{tabular}
\fi

\begin{thebibliography}{99}
\small
\raggedright
\setlength{\itemsep}{0pt}
\bibitem{erdos1958} P.~Erd\H{o}s,
  \href{https://users.renyi.hu/~p_erdos/1958-19.pdf}{\emph{Sur certaines
  s\'eries \`a valeur irrationnelle}},
  Enseign. Math. (2) \textbf{4} (1958), 93--100,
  doi:\href{https://doi.org/10.5169/seals-34629}{10.5169/seals-34629}.
  The dyadic prime series is stated as unproved on p.~94; the
  factorial-prime family is stated on p.~93, with only the $k=1$ proof printed
  on pp.~94--95.
\bibitem{erdosgraham1980} P.~Erd\H{o}s and R.~L. Graham,
  \href{https://mathweb.ucsd.edu/~ronspubs/80_11_number_theory.pdf}{\emph{Old
  and New Problems and Results in Combinatorial Number Theory}},
  Monogr. Enseign. Math. 28, Geneva, 1980, p.~62.
\bibitem{erdospomerance1978} P.~Erd\H{o}s and C.~Pomerance,
  \href{https://doi.org/10.1007/BF01818569}{\emph{On the largest prime
  factors of $n$ and $n+1$}},
  Aequationes Math. \textbf{17} (1978), 311--321,
  doi:\href{https://doi.org/10.1007/BF01818569}{10.1007/BF01818569}.
  The unnumbered dyadic irrationality theorem and its complete proof are in
  \S7 on p.~320.
\bibitem{erdos1988} P.~Erd\H{o}s, \emph{On the irrationality of certain series:
  problems and results}, in A.~Baker (ed.), \emph{New Advances in Transcendence
  Theory}, Cambridge UP, 1988, pp.~102--109,
  doi:\href{https://doi.org/10.1017/CBO9780511897184.009}{10.1017/CBO9780511897184.009}.
\bibitem{erdosstraus1963} P.~Erd\H{o}s and E.~G. Straus, \emph{On the
  irrationality of certain Ahmes series}, J. Indian Math. Soc. (N.S.)
  \textbf{27} (1964), 129--133. MR~175848.
\bibitem{fgkmt2018} K.~Ford, B.~Green, S.~Konyagin, J.~Maynard and T.~Tao,
  \emph{Long gaps between primes}, J. Amer. Math. Soc. \textbf{31} (2018),
  65--105,
  doi:\href{https://doi.org/10.1090/jams/876}{10.1090/jams/876}.
  Theorem~1 on p.~66 gives the effective lower bound for the largest single
  consecutive-prime gap below $X$.
\bibitem{kovactao2024} V.~Kova\v{c} and T.~Tao,
  \href{https://doi.org/10.1007/s10474-025-01528-0}{\emph{On several
  irrationality problems for Ahmes series}}, Acta Math. Hungar. \textbf{175}
  (2025), 572--608.
\bibitem{lean4} L.~de Moura and S.~Ullrich,
  \href{https://doi.org/10.1007/978-3-030-79876-5_37}
  {\emph{The Lean~4 theorem prover and programming language}},
  in A.~Platzer and G.~Sutcliffe (eds.), CADE~28,
  Lecture Notes in Comput. Sci. 12699, Springer, 2021, pp.~625--635,
  doi:\href{https://doi.org/10.1007/978-3-030-79876-5_37}
  {10.1007/978-3-030-79876-5\_37}.
\bibitem{mathlib} The mathlib Community,
  \href{https://doi.org/10.1145/3372885.3373824}
  {\emph{The Lean mathematical library}},
  in CPP 2020, ACM, 2020, pp.~367--381,
  doi:\href{https://doi.org/10.1145/3372885.3373824}
  {10.1145/3372885.3373824}.  The article describes a December 2019
  Lean~3-era snapshot; the repository lock owns the current revision.
\bibitem{maynard2015} J.~Maynard,
  \href{https://doi.org/10.4007/annals.2015.181.1.7}
  {\emph{Small gaps between primes}},
  Ann. of Math. (2) \textbf{181} (2015), 383--413.
\bibitem{mv2007} H.~L. Montgomery and R.~C. Vaughan,
  \href{https://doi.org/10.1017/CBO9780511618314.008}
  {\emph{Multiplicative Number Theory I: Classical Theory}},
  Cambridge Stud. Adv. Math. 97, Cambridge UP, 2007,
  Chapter~6, pp.~168--198; Theorem~6.9, pp.~179--181, and
  Exercise~6.2.5, p.~183,
  doi:\href{https://doi.org/10.1017/CBO9780511618314.008}
  {10.1017/CBO9780511618314.008}.
\bibitem{zhang2014} Y.~Zhang,
  \href{https://doi.org/10.4007/annals.2014.179.3.7}
  {\emph{Bounded gaps between primes}},
  Ann. of Math. (2) \textbf{179} (2014), 1121--1174.
\bibitem{erdosproblems} T.~F. Bloom,
  \href{https://www.erdosproblems.com/251}{\emph{Erd\H{o}s Problem \#251}},
  \texttt{erdosproblems.com/251}, accessed 28 July 2026 (page displays
  ``last edited 28 September 2025'').  The current record labels the main
  dyadic problem open, cites \texttt{[Er58b]}, \texttt{[ErGr80, p.~62]} and
  \texttt{[Er88c, p.~103]}, and explicitly describes its status as the
  website owner's present assessment rather than a literature-completeness
  guarantee; it does not mention the 2026 counterexample in
  \cite{kovac2026} to the adjacent variable-denominator conjecture.
\bibitem{kovac2026} ChatGPT~5.4 Pro (orchestrated by V.~Kova\v{c}),
  \href{https://web.math.pmf.unizg.hr/~vjekovac/files/Erdos_problem_251.pdf}
  {\emph{On the Erd\H{o}s problem \#251}},
  unpublished note, 2026, hosted by the Department of Mathematics, University
  of Zagreb, \texttt{web.math.pmf.unizg.hr}, accessed 28 July 2026.
\bibitem{formalconjectures251} The Formal Conjectures Authors,
  \href{https://github.com/google-deepmind/formal-conjectures/blob/f776d2f2039351b00737ffcafb9d7d7666e1d9af/FormalConjectures/ErdosProblems/251.lean}
  {\emph{FormalConjectures.ErdosProblems.\texttt{251}}},
  Lean source at commit \texttt{f776d2f}, 2025, accessed 28 July 2026.
\end{thebibliography}

\end{document}
